\chapter{Extended Country Profiles}
\label{app:profiles}

This appendix reproduces detailed jurisdictional narratives moved from Chapter~\ref{ch:results} to keep the results chapter within the recommended length while preserving full provenance text.

{\footnotesize
\titlespacing*{\subsection}{0pt}{0.38em}{0.08em}
\setlength{\parskip}{0pt}

\section{Detailed Narratives}
\label{sec:app-profiles}

The following subsections summarize curated narratives for each jurisdiction included in the dataset. They should be read together with theme scores and legislation lists in the JSON source of truth.


\subsection{European Union}
\textbf{Region:} Europe.
\textbf{Maturity:} Advanced.
\textbf{Primary regulator:} European Medicines Agency (EMA), National Competent Authorities, Notified Bodies.
\textbf{Privacy law:} GDPR (General Data Protection Regulation).
\textbf{AI-specific instrument:} EU AI Act (2024).
\textbf{Device framework:} Medical Device Regulation (MDR 2017/745), In Vitro Diagnostic Regulation (IVDR 2017/746), CE Marking.

\paragraph{Approval and validation:}
Conformity assessment via Notified Bodies. Risk-based classification under MDR. EU AI Act classifies most healthcare AI as high-risk requiring conformity assessments, risk management, and human oversight.
Clinical evaluation required under MDR. Post-market clinical follow-up (PMCF). EU AI Act mandates testing with real-world representative datasets, bias monitoring, and documentation.

\paragraph{Data governance and transparency:}
GDPR provides comprehensive data protection. Right to explanation for automated decisions (Article 22). Data Protection Impact Assessments (DPIA) required for high-risk processing. European Health Data Space (EHDS) proposed for secondary use of health data.
EU AI Act mandates transparency obligations for high-risk AI: technical documentation, logging, human oversight. GDPR Article 22 provides right to meaningful information about automated decision logic.

\paragraph{Ethics, surveillance, and liability:}
EU AI Act embeds ethical principles: human oversight, non-discrimination, robustness. Ethics Guidelines for Trustworthy AI (HLEG, 2019). Fundamental Rights Impact Assessment required.
Mandatory under MDR and EU AI Act. Post-market monitoring plans, serious incident reporting, periodic safety update reports.
EU AI Liability Directive (proposed) creates presumption of causality for AI harm. Product Liability Directive revision to cover software/AI.

\paragraph{Theme scores (1--10):}
Privacy~10;
Clinical validation~9;
Approval~9;
Transparency~9;
Ethics~10;
Post-market~9;
Liability~8.
Reported AI device approvals: 600.

\paragraph{Challenges and notable developments:}
Complex multi-layered regulatory landscape. Harmonization across 27 member states. Notified Body capacity constraints. Risk of over-regulation stifling innovation.
EU AI Act is the world's first comprehensive AI law. High-risk healthcare AI must comply by 2026-2027. EHDS aims to create a unified health data framework.


\subsection{Germany}
\textbf{Region:} Europe.
\textbf{Maturity:} Advanced.
\textbf{Primary regulator:} BfArM (Federal Institute for Drugs and Medical Devices).
\textbf{Privacy law:} GDPR + BDSG (Federal Data Protection Act).
\textbf{AI-specific instrument:} EU AI Act (as EU member). National AI Strategy (2018, updated 2020)..
\textbf{Device framework:} EU MDR implemented domestically. BfArM scientific advice for AI devices..

\paragraph{Approval and validation:}
EU MDR conformity assessment pathway. BfArM provides scientific advice. DiGA (Digital Health Applications) fast-track for prescribed digital therapies covered by statutory health insurance.
DiGA requires evidence of positive healthcare effects. BfArM evaluation within 12 months. Real-world evidence accepted for DiGA provisional listing.

\paragraph{Data governance and transparency:}
GDPR + BDSG. Strict interpretation of data protection. Patient Data Protection Act (PDSG, 2021). Electronic health record (ePA) rollout. German Health Data Use Act (2024).
EU AI Act transparency requirements. BDSG automated decision provisions. Strong culture of data protection enforcement.

\paragraph{Ethics, surveillance, and liability:}
Data Ethics Commission recommendations (2019). AI Observatory at BMAS. EU AI Act ethical framework. Strong academic ethics research tradition.
EU MDR post-market surveillance. BfArM device vigilance. DiGA requires ongoing evidence generation.
Product Liability Act. EU AI Liability Directive (when enacted). Medical malpractice law.

\paragraph{Theme scores (1--10):}
Privacy~10;
Clinical validation~9;
Approval~9;
Transparency~8;
Ethics~9;
Post-market~9;
Liability~7.
Reported AI device approvals: 180.

\paragraph{Challenges and notable developments:}
Historically strict data protection culture can slow health data utilization. DiGA program growing pains. Interoperability of health IT systems.
DiGA is the world's first regulatory pathway for prescribed digital therapies reimbursed by statutory health insurance. Health Data Use Act aims to unlock research data. Leading EU voice in AI regulation.


\subsection{United States}
\textbf{Region:} North America.
\textbf{Maturity:} Advanced.
\textbf{Primary regulator:} FDA (Food and Drug Administration).
\textbf{Privacy law:} HIPAA (Health Insurance Portability and Accountability Act).
\textbf{AI-specific instrument:} FDA AI/ML-Based SaMD Action Plan (2021).
\textbf{Device framework:} FDA 510(k), De Novo, PMA pathways.

\paragraph{Approval and validation:}
Risk-based classification (Class I, II, III). AI/ML devices follow SaMD (Software as a Medical Device) framework. Predetermined Change Control Plan for adaptive algorithms.
Clinical evidence required proportional to risk. Real-World Evidence (RWE) and Real-World Data (RWD) accepted. Good Machine Learning Practice (GMLP) principles published jointly with Health Canada and UK MHRA.

\paragraph{Data governance and transparency:}
HIPAA mandates protected health information (PHI) safeguards. De-identification standards under Safe Harbor and Expert Determination. HITECH Act strengthens enforcement.
FDA encourages transparency and explainability but no binding mandate. Total Product Lifecycle (TPLC) approach for monitoring algorithm changes.

\paragraph{Ethics, surveillance, and liability:}
No federal AI ethics law. NIST AI Risk Management Framework (2023) provides voluntary guidance. Executive Order on AI (2023) directs agencies to manage AI risks.
Mandatory adverse event reporting (MedWatch). Real-world performance monitoring encouraged under TPLC.
Product liability under state tort law. FDA clearance does not shield from liability. Evolving case law on AI-related malpractice.

\paragraph{Theme scores (1--10):}
Privacy~8;
Clinical validation~9;
Approval~9;
Transparency~6;
Ethics~6;
Post-market~8;
Liability~7.
Reported AI device approvals: 950.

\paragraph{Challenges and notable developments:}
Fragmented state-level privacy laws. No comprehensive federal AI legislation. Balancing innovation speed with safety oversight for adaptive algorithms.
FDA has authorized over 950 AI/ML-enabled medical devices as of 2025. Predetermined Change Control Plan allows manufacturers to pre-specify algorithm modifications.


\subsection{United Kingdom}
\textbf{Region:} Europe.
\textbf{Maturity:} Advanced.
\textbf{Primary regulator:} MHRA (Medicines and Healthcare products Regulatory Agency).
\textbf{Privacy law:} UK GDPR \& Data Protection Act 2018.
\textbf{AI-specific instrument:} Pro-innovation AI Regulation Framework (2023).
\textbf{Device framework:} UK MDR 2002 (amended), UKCA marking.

\paragraph{Approval and validation:}
Risk-based device classification similar to EU MDR but diverging post-Brexit. MHRA Software and AI as Medical Device Change Programme. Proportionate regulation favoring innovation.
NICE Evidence Standards Framework for digital health technologies. Multi-Agency Advisory Service (MAAS) provides coordinated guidance. GMLP principles co-developed with FDA and Health Canada.

\paragraph{Data governance and transparency:}
UK GDPR mirrors EU GDPR with domestic adjustments. NHS Data Strategy promotes secure data sharing. Caldicott Principles govern patient data confidentiality.
Algorithmic Transparency Recording Standard for public sector. Pro-innovation framework uses existing regulators rather than new AI-specific rules.

\paragraph{Ethics, surveillance, and liability:}
Five cross-sector AI principles: safety, transparency, fairness, accountability, contestability. AI Safety Institute established (2023). No binding AI-specific ethics law.
MHRA vigilance and post-market surveillance obligations. Yellow Card scheme for adverse incident reporting.
Common law negligence and product liability. Government exploring AI-specific liability reforms.

\paragraph{Theme scores (1--10):}
Privacy~8;
Clinical validation~8;
Approval~7;
Transparency~7;
Ethics~7;
Post-market~7;
Liability~6.
Reported AI device approvals: 200.

\paragraph{Challenges and notable developments:}
Post-Brexit regulatory divergence from EU. Balancing light-touch innovation-friendly approach with patient safety. Transition from CE to UKCA marking.
UK positioned as pro-innovation alternative to EU AI Act. AI Safety Institute is a world-first government body for AI safety research. NHS AI Lab supports AI adoption in health services.


\subsection{Canada}
\textbf{Region:} North America.
\textbf{Maturity:} Advanced.
\textbf{Primary regulator:} Health Canada.
\textbf{Privacy law:} PIPEDA (Personal Information Protection and Electronic Documents Act), Provincial health privacy laws (e.g., Ontario PHIPA).
\textbf{AI-specific instrument:} AIDA (Artificial Intelligence and Data Act, proposed in Bill C-27).
\textbf{Device framework:} Medical Devices Regulations (SOR/98-282). Risk-based classification (Class I-IV)..

\paragraph{Approval and validation:}
Risk-based classification. Pre-market review for Class II-IV. Health Canada guidance on SaMD and machine learning. GMLP principles co-developed with FDA and MHRA.
Clinical evidence requirements proportional to risk. Acceptance of RWE. Health Canada actively engaged in international harmonization (IMDRF).

\paragraph{Data governance and transparency:}
PIPEDA governs private sector. Provincial laws (PHIPA, HIA) govern health information. Proposed Consumer Privacy Protection Act (CPPA) to modernize privacy framework.
AIDA (proposed) would require transparency measures for high-impact AI. Directive on Automated Decision-Making for government systems (Algorithmic Impact Assessment).

\paragraph{Ethics, surveillance, and liability:}
Directive on Automated Decision-Making (2019) for federal government. AIDA (proposed) includes measures against biased AI. Pan-Canadian AI Strategy promotes responsible AI.
Mandatory problem reporting. Post-market surveillance plans. Canada Vigilance Program.
Common law negligence and product liability. AIDA (proposed) would introduce penalties for reckless/harmful AI deployment.

\paragraph{Theme scores (1--10):}
Privacy~7;
Clinical validation~8;
Approval~8;
Transparency~7;
Ethics~7;
Post-market~7;
Liability~6.
Reported AI device approvals: 120.

\paragraph{Challenges and notable developments:}
Federal-provincial jurisdiction complexity. AIDA legislation delayed. Keeping pace with rapid AI advancement in healthcare.
Canada co-developed GMLP with FDA and MHRA. Pan-Canadian Health Data Strategy in development. Strong AI research ecosystem (MILA, Vector, Amii) informing policy.


\subsection{Japan}
\textbf{Region:} Asia.
\textbf{Maturity:} Advanced.
\textbf{Primary regulator:} PMDA (Pharmaceuticals and Medical Devices Agency).
\textbf{Privacy law:} APPI (Act on Protection of Personal Information, amended 2022).
\textbf{AI-specific instrument:} AI Guidelines for Business (2024), Social Principles of Human-Centric AI (2019).
\textbf{Device framework:} Pharmaceutical and Medical Device Act (PMD Act). PMDA DASH (Design, Appropriate use, Security, Health) for SaMD..

\paragraph{Approval and validation:}
Risk-based classification. IDATEN regulatory sandbox for innovative devices. PMDA Special Approval for Emergency and Priority Review pathways. Continuous learning AI framework under discussion.
Clinical performance evaluation under PMD Act. PMDA guidance on AI medical device evaluation. Acceptance of foreign clinical data under ICH guidelines.

\paragraph{Data governance and transparency:}
APPI provides comprehensive data protection with special care for medical data. Next Generation Medical Infrastructure Act enables anonymized medical data use for research. Pseudonymized data processing provisions added in 2022 amendments.
Social Principles of Human-Centric AI emphasize transparency and explainability. No binding legal mandate. AI Guidelines for Business (2024) provide sector-neutral guidance.

\paragraph{Ethics, surveillance, and liability:}
Human-centric AI principles adopted by government (2019). Council for Social Principles of Human-Centric AI. No binding AI ethics legislation.
PMDA safety monitoring and re-examination. Good Vigilance Practice (GVP). Post-market clinical studies as needed.
Product Liability Act covers defective products. AI-specific liability under discussion.

\paragraph{Theme scores (1--10):}
Privacy~7;
Clinical validation~8;
Approval~8;
Transparency~6;
Ethics~7;
Post-market~8;
Liability~5.
Reported AI device approvals: 150.

\paragraph{Challenges and notable developments:}
Aging society drives demand but conservative regulatory culture. Language barriers in international harmonization. Slow adoption of digital health relative to technological capacity.
Japan's regulatory sandbox (IDATEN) accelerates innovative device approval. Society 5.0 vision integrates AI in healthcare transformation. PMDA actively developing AI/ML device evaluation frameworks.


\subsection{Singapore}
\textbf{Region:} Asia.
\textbf{Maturity:} Advanced.
\textbf{Primary regulator:} HSA (Health Sciences Authority).
\textbf{Privacy law:} PDPA (Personal Data Protection Act, 2012, amended 2021).
\textbf{AI-specific instrument:} Model AI Governance Framework (2019, updated 2020), AI Verify (2022).
\textbf{Device framework:} Health Products Act. Risk-based classification aligned with IMDRF..

\paragraph{Approval and validation:}
Risk-based classification. HSA regulatory sandbox initiatives. Aligned with IMDRF SaMD guidance. Streamlined approval for lower-risk devices.
HSA clinical evidence requirements. Recognition of international approvals. Pragmatic approach balancing evidence rigor with access.

\paragraph{Data governance and transparency:}
PDPA governs data protection. Health Information Bill (proposed) for health-specific data governance. National Electronic Health Record (NEHR) framework.
Model AI Governance Framework provides voluntary guidance on explainability. AI Verify is a testing framework for verifying AI governance claims.

\paragraph{Ethics, surveillance, and liability:}
Model AI Governance Framework emphasizes human-centric, explainable, fair AI. Advisory Council on Ethical Use of AI and Data. Smart Nation initiative integrates responsible AI.
HSA post-market surveillance. Adverse event reporting. Periodic review of registered products.
Common law product liability. Consumer Protection (Fair Trading) Act. No AI-specific liability provisions.

\paragraph{Theme scores (1--10):}
Privacy~7;
Clinical validation~7;
Approval~7;
Transparency~8;
Ethics~8;
Post-market~6;
Liability~5.
Reported AI device approvals: 60.

\paragraph{Challenges and notable developments:}
Small domestic market limits local clinical evidence generation. Balancing business-friendly environment with robust oversight. Regional influence vs. global harmonization.
AI Verify is the world's first AI governance testing framework. Singapore positioned as AI governance thought leader in ASEAN. Strong public-private partnerships in health AI.


\subsection{Switzerland}
\textbf{Region:} Europe.
\textbf{Maturity:} Advanced.
\textbf{Primary regulator:} Swissmedic.
\textbf{Privacy law:} nFADP (New Federal Act on Data Protection, 2023).
\textbf{AI-specific instrument:} No AI-specific law. Federal Council AI guidelines. Bilateral alignment with EU standards..
\textbf{Device framework:} Medical Devices Ordinance (MedDO) aligned with EU MDR. Mutual Recognition Agreement with EU..

\paragraph{Approval and validation:}
Aligned with EU conformity assessment. Swissmedic registration. Recognition of CE marking through bilateral agreements. Swiss market access closely tied to EU regulatory decisions.
Clinical evidence requirements aligned with EU MDR. Strong clinical research infrastructure. International collaboration through IMDRF.

\paragraph{Data governance and transparency:}
nFADP (2023) modernized data protection aligned with GDPR adequacy. Health data as sensitive personal data. Cantonal health data regulations.
Federal Council guidelines on AI transparency. No binding mandate. Swiss Digital Initiative Trust Label program.

\paragraph{Ethics, surveillance, and liability:}
Federal Ethics Committee on Non-Human Biotechnology. Swiss National Science Foundation AI ethics research. No binding AI ethics law.
Swissmedic vigilance. Aligned with EU post-market surveillance requirements. Adverse event reporting.
Product Liability Act. Code of Obligations. No AI-specific liability provisions.

\paragraph{Theme scores (1--10):}
Privacy~8;
Clinical validation~8;
Approval~8;
Transparency~6;
Ethics~6;
Post-market~7;
Liability~5.
Reported AI device approvals: 90.

\paragraph{Challenges and notable developments:}
Dependency on EU regulatory decisions without EU membership voting rights. Maintaining bilateral agreements post-EU AI Act. Small market size.
Strong pharma/medtech hub (Novartis, Roche). nFADP brings Swiss privacy law to GDPR-equivalent. Active AI health research at ETH Zurich and EPFL.


\subsection{China}
\textbf{Region:} Asia.
\textbf{Maturity:} Advanced.
\textbf{Primary regulator:} NMPA (National Medical Products Administration).
\textbf{Privacy law:} PIPL (Personal Information Protection Law, 2021).
\textbf{AI-specific instrument:} Interim Measures for Generative AI (2023), Algorithm Recommendation Regulations (2022), Deep Synthesis Regulations (2023).
\textbf{Device framework:} Regulations for Supervision and Administration of Medical Devices. Three-tier classification system. AI medical devices regulated as SaMD..

\paragraph{Approval and validation:}
Three-class risk-based system. NMPA has approved AI diagnostic devices in radiology, pathology, and ophthalmology. Expedited review pathways for innovative devices. Guiding Principles for AI Medical Device Registration (2021).
Clinical trials required for Class II and III devices. Multi-center clinical validation emphasized. NMPA guidance on AI/ML algorithm validation and dataset requirements.

\paragraph{Data governance and transparency:}
PIPL governs personal information with strict consent and cross-border transfer rules. Data Security Law (2021) classifies health data as important data. Sensitive personal information protections for biometric and health data.
Algorithm filing and registration required. Algorithmic Recommendation Regulations mandate transparency and user control. Deep synthesis (deepfake) regulations require labeling.

\paragraph{Ethics, surveillance, and liability:}
New Generation AI Ethics Code (2021). National AI governance principles emphasize harmony, fairness, and responsibility. State Council AI development plans integrate ethical guardrails.
Mandatory adverse event reporting to NMPA. Periodic re-evaluation of registered devices. Quality management system audits.
Civil Code provisions on product liability. Evolving jurisprudence on AI-caused harm. Algorithmic accountability provisions in algorithm regulations.

\paragraph{Theme scores (1--10):}
Privacy~7;
Clinical validation~7;
Approval~7;
Transparency~7;
Ethics~6;
Post-market~7;
Liability~6.
Reported AI device approvals: 300.

\paragraph{Challenges and notable developments:}
Rapid regulatory evolution creates compliance complexity. Data localization requirements limit international collaboration. Balancing state surveillance priorities with individual privacy.
China is among the first to regulate algorithmic recommendations and generative AI. NMPA rapidly approving AI diagnostic tools. National AI standardization efforts underway.


\subsection{Australia}
\textbf{Region:} Oceania.
\textbf{Maturity:} Moderate.
\textbf{Primary regulator:} TGA (Therapeutic Goods Administration).
\textbf{Privacy law:} Privacy Act 1988 (amended), Australian Privacy Principles (APPs).
\textbf{AI-specific instrument:} Australia's AI Ethics Principles (2019), Voluntary AI Safety Standard (2024).
\textbf{Device framework:} Therapeutic Goods Act 1989. Risk-based classification aligned with IMDRF. TGA regulatory guidance on SaMD..

\paragraph{Approval and validation:}
Risk-based classification. TGA registration and conformity assessment. Alignment with international standards. SaMD-specific regulatory guidance.
Clinical evidence requirements. TGA guidance on digital health technologies. Recognition of international clinical evidence.

\paragraph{Data governance and transparency:}
Privacy Act with Australian Privacy Principles. My Health Records Act for digital health records. Health data subject to enhanced protections.
Voluntary AI Ethics Principles include transparency. No binding mandate. Government response to AI consultation ongoing.

\paragraph{Ethics, surveillance, and liability:}
Eight AI Ethics Principles (2019): human-centric, fair, transparent, contestable, accountable, reliable, safe, privacy-protected. Voluntary adoption. National AI Centre provides guidance.
TGA post-market review and vigilance. Adverse event reporting mandatory. Systematic review of registered products.
Australian Consumer Law product liability. Professional negligence standards. No AI-specific liability law.

\paragraph{Theme scores (1--10):}
Privacy~7;
Clinical validation~7;
Approval~7;
Transparency~6;
Ethics~7;
Post-market~7;
Liability~5.
Reported AI device approvals: 80.

\paragraph{Challenges and notable developments:}
Voluntary nature of AI ethics framework limits enforcement. Resource constraints in TGA for AI evaluation. Geographic barriers to clinical trial recruitment.
Australia's AI Ethics Principles among the earliest globally. TGA actively developing SaMD regulatory pathways. National Digital Health Strategy driving AI integration.


\subsection{South Korea}
\textbf{Region:} Asia.
\textbf{Maturity:} Advanced.
\textbf{Primary regulator:} MFDS (Ministry of Food and Drug Safety).
\textbf{Privacy law:} PIPA (Personal Information Protection Act, amended 2023).
\textbf{AI-specific instrument:} AI Framework Act (proposed), National AI Strategy (2019).
\textbf{Device framework:} Medical Devices Act. Risk-based classification. Regulatory Science Center for AI/ML devices..

\paragraph{Approval and validation:}
Risk-based classification. MFDS has fast-track pathways for innovative AI devices. Regulatory Science Center specializing in AI medical device review. Software-specific classification criteria.
Clinical trial requirements under Medical Devices Act. MFDS guidance on AI/ML medical device clinical validation. Growing emphasis on real-world evidence.

\paragraph{Data governance and transparency:}
PIPA provides comprehensive data protection. 2023 amendments enable pseudonymized data processing for research. MyData framework for personal data portability.
National AI Strategy emphasizes trustworthy AI. No binding transparency mandate yet. AI Framework Act (proposed) would establish transparency requirements.

\paragraph{Ethics, surveillance, and liability:}
National AI Ethics Standards (2020). AI Ethics Guidelines emphasize human dignity, public good, and technical soundness. No binding legislation yet.
MFDS post-market surveillance. Adverse event reporting system. Quality management requirements.
Product Liability Act. Civil liability framework. AI-specific liability considerations under development.

\paragraph{Theme scores (1--10):}
Privacy~7;
Clinical validation~7;
Approval~8;
Transparency~6;
Ethics~6;
Post-market~7;
Liability~5.
Reported AI device approvals: 200.

\paragraph{Challenges and notable developments:}
Balancing rapid technological adoption with regulatory oversight. Harmonizing domestic standards with international frameworks. SME compliance burden.
South Korea has one of the highest AI medical device approval rates globally. MFDS Regulatory Science Center is a dedicated AI device evaluation unit. K-BioHealth initiative driving AI healthcare innovation.


\subsection{Israel}
\textbf{Region:} Middle East.
\textbf{Maturity:} Advanced.
\textbf{Primary regulator:} MOH AMAR (Division of Medical Devices, Ministry of Health).
\textbf{Privacy law:} Privacy Protection Law (1981, amended), Health Information Exchange regulations.
\textbf{AI-specific instrument:} AI Policy and Ethics Framework (2022), National AI Program.
\textbf{Device framework:} Medical Device Regulations under Pharmacists' Ordinance. Alignment with EU MDR and FDA pathways..

\paragraph{Approval and validation:}
Recognized market approval from FDA, CE, or other reference agencies. AMAR reviews and registers. Expedited pathways for innovative technologies. Growing AI device registrations.
Clinical trial requirements under MOH. Strong real-world data infrastructure. International clinical validation partnerships.

\paragraph{Data governance and transparency:}
Privacy Protection Law with sector-specific health regulations. Unique national health data infrastructure (Clalit, Maccabi, etc.). Rich longitudinal health datasets enabling AI research.
AI Policy Framework recommends transparency. No binding mandate. Active academic and industry discourse on AI governance.

\paragraph{Ethics, surveillance, and liability:}
AI Ethics Framework addresses healthcare applications. Innovation Authority promotes responsible AI. Active ethics research community.
MOH post-market surveillance. Adverse event reporting. Device vigilance system.
Tort liability. Patient Rights Law. No AI-specific liability provisions. Active legal scholarship on the topic.

\paragraph{Theme scores (1--10):}
Privacy~6;
Clinical validation~8;
Approval~7;
Transparency~6;
Ethics~6;
Post-market~6;
Liability~5.
Reported AI device approvals: 100.

\paragraph{Challenges and notable developments:}
Outdated base privacy legislation. Balancing startup-friendly ecosystem with regulatory rigor. Small domestic market necessitates international focus.
Israel is a global health AI innovation hub. Unique national health data infrastructure enables world-class AI research. Highest health AI startups per capita globally.


\subsection{Brazil}
\textbf{Region:} South America.
\textbf{Maturity:} Developing.
\textbf{Primary regulator:} ANVISA (Agência Nacional de Vigilância Sanitária).
\textbf{Privacy law:} LGPD (Lei Geral de Proteção de Dados, 2020).
\textbf{AI-specific instrument:} AI Bill (PL 2338/2023, under legislative process).
\textbf{Device framework:} RDC 185/2001 and updates. Risk-based classification (Classes I-IV)..

\paragraph{Approval and validation:}
Risk-based classification. ANVISA registration required. SaMD guidance under development. International harmonization through IMDRF participation.
Clinical evidence requirements under ANVISA regulations. Limited AI-specific validation guidance. Growing clinical AI research community.

\paragraph{Data governance and transparency:}
LGPD provides comprehensive data protection modeled on GDPR. Health data classified as sensitive. ANPD (National Data Protection Authority) oversees enforcement.
AI Bill (proposed) includes transparency obligations for high-risk AI. LGPD provides right to review of automated decisions. No current binding AI transparency mandate.

\paragraph{Ethics, surveillance, and liability:}
AI Bill (proposed) incorporates ethical principles: human oversight, non-discrimination, transparency. Brazilian AI Strategy (2021) promotes responsible development.
ANVISA post-market surveillance. Tecnovigilância (technovigilance) system. Mandatory adverse event reporting.
Consumer Defense Code. Civil liability framework. AI Bill (proposed) would establish risk-based liability for AI.

\paragraph{Theme scores (1--10):}
Privacy~7;
Clinical validation~5;
Approval~5;
Transparency~5;
Ethics~5;
Post-market~5;
Liability~5.
Reported AI device approvals: 40.

\paragraph{Challenges and notable developments:}
AI regulatory framework still in legislative process. Limited ANVISA capacity for AI device evaluation. Digital health infrastructure gaps in underserved regions.
LGPD is Latin America's most comprehensive data protection law. Active AI legislation process. Growing health AI startup ecosystem in São Paulo.


\subsection{UAE}
\textbf{Region:} Middle East.
\textbf{Maturity:} Emerging.
\textbf{Primary regulator:} MOH (Ministry of Health and Prevention), HAAD, DHA.
\textbf{Privacy law:} Federal Decree-Law No. 45/2021 (Personal Data Protection).
\textbf{AI-specific instrument:} National AI Strategy 2031, Minister of State for AI appointed (2017).
\textbf{Device framework:} MOH medical device registration. Emirates Conformity Assessment Scheme (ECAS)..

\paragraph{Approval and validation:}
MOH registration with recognition of international approvals (FDA, CE). Emirates Authority for Standardization (ESMA) standards. Dubai Health Authority (DHA) parallel pathway. Expedited processes for innovative technologies.
Clinical evidence requirements for device registration. Recognition of international clinical data. Mohamed bin Rashid University of Medicine and Health Sciences research programs.

\paragraph{Data governance and transparency:}
Federal data protection law (2021). Dubai International Financial Centre (DIFC) Data Protection Law. Abu Dhabi Healthcare Information and Cyber Security Standard. Dubai Health Data Policy.
National AI Strategy includes governance principles. No binding transparency mandate. AI governance guidelines under development.

\paragraph{Ethics, surveillance, and liability:}
National AI Strategy ethical guidelines. UAE AI Ethics Board. Oxford-UAE AI research partnership on responsible AI.
MOH device vigilance. Adverse event reporting. Quality management requirements.
Civil Code liability provisions. Consumer Protection Law. No AI-specific liability framework.

\paragraph{Theme scores (1--10):}
Privacy~5;
Clinical validation~5;
Approval~5;
Transparency~5;
Ethics~5;
Post-market~4;
Liability~3.
Reported AI device approvals: 35.

\paragraph{Challenges and notable developments:}
Multi-authority regulatory landscape (federal and emirate-level). Rapid strategy ambitions vs. regulatory maturity. Building domestic healthcare data infrastructure.
First country to appoint a Minister of State for AI (2017). Ambitious National AI Strategy 2031. Hub for health AI startups and international collaborations.


\subsection{Saudi Arabia}
\textbf{Region:} Middle East.
\textbf{Maturity:} Emerging.
\textbf{Primary regulator:} SFDA (Saudi Food and Drug Authority).
\textbf{Privacy law:} PDPL (Personal Data Protection Law, 2023).
\textbf{AI-specific instrument:} SDAIA AI Ethics Principles (2023), National AI Strategy.
\textbf{Device framework:} Medical Devices Interim Regulation. SFDA SaMD guidance. Alignment with IMDRF..

\paragraph{Approval and validation:}
SFDA registration and conformity assessment. Recognition of international approvals (FDA, CE). Expedited review pathways for innovative products. SaMD regulatory framework under development.
Clinical evidence requirements. Recognition of international clinical data. National Center for AI (NCAI) supports validation research.

\paragraph{Data governance and transparency:}
PDPL (effective 2023) establishes comprehensive data protection. Health data classified as sensitive. Cloud Computing Regulatory Framework governs data hosting.
SDAIA AI Ethics Principles include transparency. No binding mandate. Data and AI governance frameworks under development.

\paragraph{Ethics, surveillance, and liability:}
SDAIA AI Ethics Principles (2023): fairness, transparency, accountability, privacy, security, human-centricity. Vision 2030 AI integration goals.
SFDA post-market surveillance requirements. Adverse event reporting. Medical device vigilance system.
Civil liability framework. AI-specific liability not yet addressed. Consumer protection provisions apply.

\paragraph{Theme scores (1--10):}
Privacy~5;
Clinical validation~4;
Approval~5;
Transparency~5;
Ethics~5;
Post-market~4;
Liability~3.
Reported AI device approvals: 25.

\paragraph{Challenges and notable developments:}
Regulatory framework still maturing. Building domestic regulatory expertise. Dependency on international recognition pathways. Limited public health data infrastructure.
Massive AI investment under Vision 2030. SDAIA established as dedicated AI authority. NEOM project integrating AI healthcare. Rapid digital health transformation.


\subsection{India}
\textbf{Region:} Asia.
\textbf{Maturity:} Developing.
\textbf{Primary regulator:} CDSCO (Central Drugs Standard Control Organisation).
\textbf{Privacy law:} Digital Personal Data Protection Act (DPDP, 2023).
\textbf{AI-specific instrument:} No AI-specific law. NITI Aayog AI Strategy and Responsible AI guidelines (2021)..
\textbf{Device framework:} Medical Devices Rules 2017 (amended 2020). SaMD classification under development..

\paragraph{Approval and validation:}
Risk-based classification (Class A-D). CDSCO registration required. Limited specific guidance for AI/ML medical devices. Reliance on international standards (ISO, IEC).
Clinical investigation requirements under Medical Devices Rules. Limited AI-specific validation guidance. Growing clinical AI research ecosystem.

\paragraph{Data governance and transparency:}
DPDP Act 2023 establishes consent-based framework with data fiduciary obligations. Health data treated as sensitive. Ayushman Bharat Digital Mission (ABDM) creates digital health infrastructure.
No binding transparency requirements. NITI Aayog principles recommend explainability and accountability. IndiaAI Mission (2024) promotes responsible AI.

\paragraph{Ethics, surveillance, and liability:}
NITI Aayog Responsible AI principles (fairness, accountability, transparency, privacy, safety). No binding ethics law for AI. ICMR National Ethical Guidelines for Biomedical and Health Research.
Post-market surveillance requirements under Medical Devices Rules. Materiovigilance Programme of India (MvPI) for adverse event reporting.
Consumer Protection Act 2019 covers product defects. IT Act provisions for data breaches. No AI-specific liability framework.

\paragraph{Theme scores (1--10):}
Privacy~5;
Clinical validation~4;
Approval~4;
Transparency~4;
Ethics~5;
Post-market~4;
Liability~4.
Reported AI device approvals: 30.

\paragraph{Challenges and notable developments:}
Regulatory framework still maturing for AI devices. Limited CDSCO capacity for AI/ML evaluation. Infrastructure gaps in digital health. Balancing innovation needs with regulatory rigor.
DPDP Act 2023 is India's first comprehensive data protection law. Ayushman Bharat Digital Mission building foundational digital health infrastructure. IndiaAI Mission allocating significant resources for AI development.


\subsection{Thailand}
\textbf{Region:} Asia.
\textbf{Maturity:} Developing.
\textbf{Primary regulator:} Thai FDA (Food and Drug Administration).
\textbf{Privacy law:} PDPA (Personal Data Protection Act, 2022).
\textbf{AI-specific instrument:} National AI Strategy and Action Plan (2022-2027). No binding AI law..
\textbf{Device framework:} Medical Device Act B.E. 2551 (2008). Risk-based classification. SaMD guidance under development..

\paragraph{Approval and validation:}
Thai FDA registration. Risk-based classification. Limited AI-specific device guidance. ASEAN Medical Device Directive harmonization.
Clinical evidence requirements. Thai Clinical Trials Registry. Limited AI-specific validation frameworks. Growing research capacity.

\paragraph{Data governance and transparency:}
PDPA provides comprehensive data protection effective 2022. Health data as sensitive data. Ministry of Public Health data governance policies.
National AI Ethics Guidelines include transparency. No binding mandate. NECTEC AI research addresses explainability.

\paragraph{Ethics, surveillance, and liability:}
National AI Ethics Guidelines. Thailand 4.0 policy includes responsible technology principles. No binding AI ethics law.
Thai FDA post-market surveillance. Adverse event reporting. Resource-constrained monitoring.
Consumer Protection Act. Product Liability Act. No AI-specific liability.

\paragraph{Theme scores (1--10):}
Privacy~5;
Clinical validation~4;
Approval~4;
Transparency~4;
Ethics~4;
Post-market~3;
Liability~3.
Reported AI device approvals: 15.

\paragraph{Challenges and notable developments:}
Regulatory framework for AI devices still developing. Digital health infrastructure building. Balancing medical tourism innovation with domestic healthcare priorities.
PDPA implementation is a major privacy milestone. Thailand 4.0 drives digital health investment. Medical tourism sector catalyzing AI adoption.


\subsection{South Africa}
\textbf{Region:} Africa.
\textbf{Maturity:} Emerging.
\textbf{Primary regulator:} SAHPRA (South African Health Products Regulatory Authority).
\textbf{Privacy law:} POPIA (Protection of Personal Information Act, 2021).
\textbf{AI-specific instrument:} No AI-specific regulation. AI Policy Framework under development..
\textbf{Device framework:} Medicines and Related Substances Act. Medical device regulations under SAHPRA. Classification system under implementation..

\paragraph{Approval and validation:}
SAHPRA registration transitioning from previous MCC system. Medical device regulatory framework being modernized. Limited AI-specific device guidance. Backlog in device registrations.
Clinical evidence requirements under development. Reliance on international clinical data. Limited local AI validation infrastructure.

\paragraph{Data governance and transparency:}
POPIA provides comprehensive data protection. Health data classified as special personal information. Conditions for lawful processing include consent and public health exceptions.
No binding transparency requirements. Presidential Commission on 4IR (2020) recommended AI governance. Draft AI Policy Framework under consultation.

\paragraph{Ethics, surveillance, and liability:}
Presidential Commission on 4IR recommendations. No binding AI ethics framework. Research ethics committees govern clinical AI research. Ubuntu philosophy informing ethical discourse.
SAHPRA post-market surveillance developing. Adverse event reporting being strengthened. Resource limitations affect monitoring capacity.
Consumer Protection Act. Common law product liability. No AI-specific liability provisions.

\paragraph{Theme scores (1--10):}
Privacy~6;
Clinical validation~3;
Approval~3;
Transparency~3;
Ethics~4;
Post-market~3;
Liability~3.
Reported AI device approvals: 10.

\paragraph{Challenges and notable developments:}
SAHPRA capacity and resource constraints. Digital infrastructure gaps. Brain drain in health AI expertise. Balancing innovation with addressing basic healthcare access.
POPIA is Africa's most comprehensive data protection law. CSIR and universities driving AI health research. Growing health tech ecosystem in Cape Town and Johannesburg.


\subsection{Kenya}
\textbf{Region:} Africa.
\textbf{Maturity:} Early.
\textbf{Primary regulator:} PPB (Pharmacy and Poisons Board).
\textbf{Privacy law:} Data Protection Act (2019).
\textbf{AI-specific instrument:} No AI-specific regulation. Blockchain and AI Taskforce Report (2019)..
\textbf{Device framework:} PPB medical device regulations. Classification aligned with WHO guidelines. Limited SaMD framework..

\paragraph{Approval and validation:}
PPB registration for medical devices. Limited AI-specific guidance. Recognition of WHO prequalified products. Growing regulatory capacity through international partnerships.
National Commission for Science, Technology and Innovation (NACOSTI) oversees research. Clinical trial regulations. Limited AI-specific validation frameworks.

\paragraph{Data governance and transparency:}
Data Protection Act 2019 with Office of Data Protection Commissioner (ODPC). Health data as sensitive data category. Kenya Health Policy framework guides health data governance.
No binding transparency requirements. Blockchain and AI Taskforce recommendations (2019) address governance. Limited implementation.

\paragraph{Ethics, surveillance, and liability:}
Blockchain and AI Taskforce ethical recommendations. No binding AI ethics framework. KEMRI Ethics Review Committee for health research.
PPB pharmacovigilance. Limited medical device post-market monitoring. Resource constraints.
Consumer Protection Act. Tort law framework. No AI-specific liability provisions.

\paragraph{Theme scores (1--10):}
Privacy~5;
Clinical validation~2;
Approval~3;
Transparency~2;
Ethics~3;
Post-market~2;
Liability~2.
Reported AI device approvals: 3.

\paragraph{Challenges and notable developments:}
Limited regulatory capacity. Infrastructure constraints. Prioritizing basic healthcare access alongside innovation. Dependency on international standards and partnerships.
Kenya is a mobile health innovation leader in Africa. Data Protection Act among earliest in East Africa. Growing AI-for-health research at universities and tech hubs.


\subsection{Nigeria}
\textbf{Region:} Africa.
\textbf{Maturity:} Early.
\textbf{Primary regulator:} NAFDAC (National Agency for Food and Drug Administration and Control).
\textbf{Privacy law:} NDPR (Nigeria Data Protection Regulation, 2019) / NDPA (Nigeria Data Protection Act, 2023).
\textbf{AI-specific instrument:} National AI Strategy (draft). No binding AI regulation..
\textbf{Device framework:} NAFDAC medical device registration. Classification system. Limited SaMD-specific framework..

\paragraph{Approval and validation:}
NAFDAC registration for medical devices. Limited specific guidance for AI/ML devices. Reliance on international approvals and standards.
Limited AI-specific clinical validation requirements. Research ethics committees govern clinical studies. Growing AI-for-health research community.

\paragraph{Data governance and transparency:}
NDPA 2023 establishes comprehensive data protection. Health data as sensitive category. NITDA (National Information Technology Development Agency) oversees digital policy.
No binding transparency requirements. National Digital Economy Policy addresses some digital governance aspects.

\paragraph{Ethics, surveillance, and liability:}
No binding AI ethics framework. Academic and civil society discourse on AI ethics. NITDA strategic framework includes responsible AI elements.
NAFDAC pharmacovigilance. Limited medical device post-market monitoring. Resource constraints affect surveillance capacity.
Consumer protection laws. Federal Competition and Consumer Protection Act. No AI-specific liability.

\paragraph{Theme scores (1--10):}
Privacy~4;
Clinical validation~2;
Approval~3;
Transparency~2;
Ethics~3;
Post-market~2;
Liability~2.
Reported AI device approvals: 5.

\paragraph{Challenges and notable developments:}
Limited regulatory capacity for AI devices. Digital infrastructure gaps. Healthcare access challenges precede AI adoption. Funding constraints for regulatory modernization.
NDPA 2023 is a significant step in data governance. Growing AI research community (e.g., Data Science Nigeria). Mobile health innovations bypassing traditional regulatory pathways.

} % end compact appendix body


